\section{Variables que solo registra un canal: el fallo visible}\label{sec:canal}

Una covariable que solo se registra en uno de los canales rompe el soporte aunque las unidades sean comparables. Si no predice el resultado dadas las demás, descartarla restaura la identificación; si confunde, ni usarla ni descartarla identifica $\theta$.

\paragraph{Modelo de medición.} Cada unidad tiene covariables latentes $X^*=(V,C^*)$, donde $C^*$ toma valores en $\mathcal C\cup\{\vacio\}$ y $C^*=\vacio$ significa «sin dato real». Lo registrado depende del canal: $X=(V,C)$, con $C=C^*$ si $D=1$ y $C=\vacio$ si $D=0$. La sección~\ref{sec:instrumentos} trata el caso general $X=\varphi_D(X^*)$. Ejemplos: un historial comercial que solo guardan los gestores del canal premium, o pruebas diagnósticas que solo se hacen en el hospital. Escribimos $P^*$ y $Q^*$ para las leyes de $X^*$ en el grupo objetivo y en el canal donante.

\begin{lema}[Masa singular]\label{lem:singular}
$\varepsilon=\Ps(\X)=\sup\{P(B):\ B\in\A,\ Q(B)=0\}$.
\end{lema}
\begin{proof}
Sea $N$ un conjunto $Q$-nulo donde se concentra $\Ps$. Para todo $B$ con $Q(B)=0$, $\Pac(B)=0$, porque $\Pac\ll Q$; luego $P(B)=\Ps(B)\le\Ps(\X)$. Además $P(N)=\Pac(N)+\Ps(N)=\Ps(\X)$, así que el supremo se alcanza en $N$.
\end{proof}

\begin{proposicion}[Una variable de canal rompe el soporte]\label{prop:rompe}
Bajo el modelo de medición, sea $\kappa=\PP(C^*\ne\vacio\mid S=1)$ la fracción del grupo objetivo con dato real. Entonces $\varepsilon\ge\kappa$, y el clasificador ideal vale 1 $M$-casi seguro en los perfiles con $C\ne\vacio$. Si $\kappa>0$, (A3) falla con independencia de cómo sean las unidades.
\end{proposicion}
\begin{proof}
Sea $B=\mathcal V\times\mathcal C$, el conjunto de perfiles con dato registrado. Por el modelo de medición, $Q(B)=\PP(C\ne\vacio\mid D=0)=0$, mientras que $P(B)=\kappa$; el Lema~\ref{lem:singular} da $\varepsilon\ge\kappa$. Como $\PP(Z=0,\tilde X\in B)=(1-\pi)Q(B)=0$, $p=1$ $M$-casi seguro sobre $B$.
\end{proof}

En adelante, $g$ denota una versión de $\E[Y(0)\mid V,D=0]$, y $P_V$ y $Q_V$ las leyes de $V$ en el grupo objetivo y en el canal donante.

\begin{proposicion}[Descartar una variable que no predice el resultado]\label{prop:descartar}
Supongamos:
\begin{enumerate}
  \item \emph{Soporte latente:} $P^*\ll Q^*$.
  \item \emph{(A2) latente:} $\E[Y(0)\mid X^*,S=1]=\E[Y(0)\mid X^*,D=0]$, $P^*$-casi seguro.
  \item \emph{$C^*$ no predice el resultado en el canal básico dado $V$:} $\E[Y(0)\mid X^*,D=0]=\E[Y(0)\mid V,D=0]$, $Q^*$-casi seguro.
\end{enumerate}
Entonces, con $V$ como únicas covariables, (a) $P_V\ll Q_V$ y (b) $\E[Y(0)\mid V,S=1]=\E[Y(0)\mid V,D=0]$ $P_V$-casi seguro, así que los Teoremas~\ref{teo:identificacion}--\ref{teo:lebesgue} se aplican con $X=V$.
\end{proposicion}
\begin{proof}
(a) Si $Q_V(A)=0$, entonces $Q^*(A\times(\mathcal C\cup\{\vacio\}))=0$, y por (1) también $P^*(A\times(\mathcal C\cup\{\vacio\}))=P_V(A)=0$.
(b) Por (3), $\E[Y(0)\mid X^*,D=0]=g(V)$ salvo en un conjunto $Q^*$-nulo, que por (1) es también $P^*$-nulo. Por (2), $\E[Y(0)\mid X^*,S=1]=g(V)$ $P^*$-casi seguro, y como $V$ es función de $X^*$, la torre dentro de $\{S=1\}$ da $\E[Y(0)\mid V,S=1]=g(V)$.
\end{proof}

La condición (3) es un criterio de selección de covariables por el lado del resultado \citep{deluna2011,vanderweele2011}: basta aunque $C^*$ esté asociada al canal.

\paragraph{Cuando la variable sí confunde.} Si se mantienen (1) y (2) pero falla (3), sea $\mu^*(x^*)=\E[Y(0)\mid X^*=x^*,D=0]$; entonces $\theta=\int\mu^*\dd P^*$. Usando solo $V$, el estimador converge a $\int g\dd P_V=\int g(V)\dd P^*$, y su sesgo es
\[
\int g\dd P_V-\theta=-\int_{\X^*}(\mu^*-g)\dd P^*.
\]
Bajo $Q^*$ esa integral es cero por definición de $g$. Bajo $P^*$ solo puede ser distinta de cero si $C^*$ predice $Y(0)$ dado $V$ y además su distribución dado $V$ difiere entre grupos: son las dos condiciones de la definición clásica de confusor, necesarias pero no suficientes, porque los efectos podrían compensarse. La condición exacta para (A2) con $V$ como covariables es $\E_{P^*}[\mu^*-g\mid V]=0$ $P_V$-casi seguro; el sesgo en sí se anula si y solo si $\int(\mu^*-g)\dd P^*=0$, lo que también puede ocurrir por compensación. (3) es suficiente para ambas y no depende de $P^*$. Cuando el sesgo no es nulo, ninguna de las dos opciones identifica $\theta$: sin $C$ falla (A2), y con $C$ falla (A3) y, como se ve a continuación, en general también (A2) en la parte con soporte. Por eso tampoco sirven las cotas del Corolario~\ref{cor:cotas}, que suponen (A2) en esa parte. Lo que queda es un análisis de sensibilidad sobre (3).

\paragraph{Conservar la variable tampoco arregla la parte con soporte.} Con $C$ entre las covariables, la parte con soporte es $\{C=\vacio\}$. Pero «$\vacio$» no significa lo mismo en los dos grupos: en el canal donante mezcla unidades con $C^*=\vacio$ y con $C^*\ne\vacio$, porque ahí nunca se registra; en el grupo objetivo equivale a $C^*=\vacio$. Como todo el canal donante tiene $C=\vacio$, $m_0(v,\vacio)=g(v)$, y el estimador ponderado converge a $\E_{P^*}[g(V)\mid C^*=\vacio]$, que en general no es ni $\theta$ ni $\thac$. En efecto, (A2) en $\{C=\vacio\}$ exige $\mu^*(V,\vacio)=g(V)$ $P^*$-casi seguro sobre $\{C^*=\vacio\}$, es decir, que entre los donantes con el mismo $V$ respondan igual en el canal básico los que tienen historial y los que no; (3) lo implica, pero es más débil. El doblemente robusto, además, extrapola el modelo de resultado a las unidades con dato (Corolario~\ref{cor:drsinsoporte}). Ese modelo se entrena en el canal donante, donde $C$ nunca varía; si no depende de $C$ fuera del soporte, como ocurre con árboles, que no pueden partir por una variable constante en el entrenamiento, extrapola $g(V)$, y el doblemente robusto converge a $\E_{P^*}[g(V)]$, como si $C$ no se usara.

\paragraph{Contrastabilidad.} La condición (3) concierne a $Y(0)$ en el canal básico, donde $C^*$ no se observa, así que no es contrastable directamente. Que $C$ prediga el resultado en el canal premium es un indicio pero no una prueba, porque ahí se observa $Y(1)$.

\paragraph{Relación con el enfoque por patrón de ausencia.} \citet{dagostino2000} y \citet{blake2020} estiman la propensión dentro de cada patrón de valores ausentes y exigen que la variable parcialmente observada no confunda en el patrón donde falta. Una variable de canal es el caso en que el patrón lo determina el tratamiento: en el patrón observado no hay unidades donantes (Proposición~\ref{prop:rompe}), y en el ausente exige $\mu^*(V,\vacio)=g(V)$, que (3) implica.

\paragraph{Diagnóstico práctico.} Retirar una variable de canal debe reducir de forma visible el AUC del clasificador y la masa estimada $\hat P(\hat p>1-t)$; un salto grande en esa masa al incluir una variable es la firma de la Proposición~\ref{prop:rompe}. En la simulación, la masa sin soporte estimada con la variable de canal es $0,326$ frente a $\kappa=0,335$, y el estimador doblemente robusto se comporta como predice el Corolario~\ref{cor:drsinsoporte} y no como el ponderado (sección~\ref{sec:res-principal}).

\subsection{Ilustración: una referencia experimental con una covariable censurada}\label{sec:lalonde}

Una covariable no necesita faltar para romper el soporte: basta con que un canal la registre con un tope. Unos datos públicos con referencia experimental muestran un fallo así, y que el clasificador estimado no lo detecta.

\paragraph{Datos.} Usamos la submuestra de \citet{dehejia1999} del experimento National Supported Work, con 185 participantes y 260 controles aleatorizados, y el grupo de comparación CPS-1 de \citet{lalonde1986}, con 15\,992 personas. El grupo objetivo son los participantes, y los donantes, la muestra de la CPS. El estimando es la media de sus ingresos de 1978 sin el programa, cuya referencia experimental es la media de los controles: 4\,555 dólares (error típico 340). Las covariables son edad, años de estudio, raza, estado civil, falta de título, ingresos de 1974 y 1975 con indicadores de ingresos nulos, y los cuadrados de la edad y los años de estudio. Este análisis no se registró de antemano; damos la mediana sobre cinco particiones del \emph{cross-fitting}.

\paragraph{Resultados.} La media de la CPS sin ajustar es 14\,847 dólares. La ponderación sola da 5\,092 con un clasificador logístico y 3\,990 con LightGBM \citep{ke2017}, mientras que el estimador doblemente robusto da 4\,783 y 4\,716 (errores típicos 397 y 876), ambos a menos de un error típico de la referencia (340). El solapamiento es malo: con LightGBM, $\neff$ es 33 de 15\,992 donantes.

\paragraph{Tope en una de las fuentes.} Las dos fuentes registran los ingresos de forma distinta \citep{smith2005}. Una diferencia es visible: los ingresos de la CPS están censurados por arriba (el 19\,\% de los ingresos de 1974 está exactamente en el tope de 25\,862 dólares) y los del NSW no. Dos participantes (1,1\,\%) ganaron más que el tope en 1974, así que sus perfiles no tienen soporte por construcción: una masa singular en el sentido del Lema~\ref{lem:singular}, donde el clasificador ideal vale 1. Ninguno de los dos clasificadores estimados se acerca a 1: ambos dan a esas dos unidades un $\hat p$ inferior a 0,3, porque el 19\,\% de donantes en el tope hace que los ingresos altos parezcan propios de los donantes. Una falta de soporte en el extremo de una variable censurada apunta, así, el diagnóstico en la dirección equivocada. Aplicar el tope a los dos grupos cambia las estimaciones menos de 30 dólares, así que aquí no importa. El resultado también está censurado: el 26\,\% de los ingresos de 1978 de la CPS está en el tope, lo que sesga a la baja el modelo de resultado de los donantes; aplicar el tope a los controles baja la referencia en 54 dólares, a 4\,501, poco frente a los errores típicos. El tope es un instrumento propio de un canal que solo rompe el soporte en su extremo. La sección~\ref{sec:instrumentos} pasa a instrumentos que no lo rompen en absoluto; otras diferencias en cómo recogieron los ingresos las dos fuentes serían de ese tipo.
